\slajdid{04-s019}
\begin{frame}[label=04-s019]{Grupisanje celobrojnih cifara}
% element: 04-s019:e01 — Primer 1AC3_16: prevod svake cifre i uklanjanje tri vodeće nule
% element: 04-s019:e02 — Primer 10100010010101101_2: grupisanje zdesna i dodavanje nula
% element: 04-s019:e03 — Konačni rezultat 144AD_16 i analogija sa trocifrenim oktalnim grupama
\setlength{\abovedisplayskip}{3pt}\setlength{\belowdisplayskip}{3pt}
\begin{columns}[T,onlytextwidth]
\column{.475\textwidth}\NaslovBloka{Heksadecimalni $\to$ binarni}
Svaku cifru prevodimo zasebno.
\[1AC3_{16}={?}_2\]
\[=0001\ 1010\ 1100\ 0011_2\]
Brišemo vodeće nule:
\[=\alert{000}1\ 1010\ 1100\ 0011_2\]
\[\alert{1AC3_{16}=1101011000011_2}\]
\column{.475\textwidth}\NaslovBloka{Binarni $\to$ heksadecimalni}
Grupe od četiri cifre, zdesna nalevo.
\[10100010010101101_2={?}_{16}\]
\[1\ 0100\ 0100\ 1010\ 1101_2\]
Dodajemo vodeće nule:
\[\alert{000}1\ 0100\ 0100\ 1010\ 1101_2\]
\[\alert{=144AD_{16}}\]
\end{columns}
\par\vspace{3mm}\Rezultat{Za oktalni sistem postupak je isti, sa grupama od tri cifre.}
\end{frame}
